\hyperref[sec:14-Deep-Learning/07-Generative-Models-Unified/01]{``四条路线在最小化哪个散度''}把生成模型摆在同一张桌子上比过一轮，结论里有一条对这一章特别关键：真正难的从来不是写下目标，而是那个目标要求你一次性把一个极复杂的分布对齐到数据。VAE 用一个下界绕开它，GAN 用一个博弈绕开它，流模型用可逆性硬扛它，各有各的代价。

扩散模型换了一个思路，而且换得相当彻底：\textbf{不做那一次难的匹配，把它拆成很多次容易的匹配。}在数据和纯噪声之间铺一条长链，链上相邻两点近到它们之间的关系可以用一个高斯写完，于是每一小步的分布匹配都有闭式、都能训、都只需要网络前向一次。整条链走通，难的那件事就自动完成了。这个拆法带来的好处和账单都很干净：训练变成一个稳定得不像生成模型的回归问题，采样变成一道必须一步一步付钱的数值积分题。

四篇分别回答一个问题。

\subsection{本章篇目}

\begin{itemize}
\tightlist
\item \hyperref[sec:14-Deep-Learning/08-Diffusion-Models/01]{``前向加噪是一条不用学的固定路径''}：噪声怎么加。为什么单步系数是 \(\sqrt{1-\beta_t}\)，为什么整条链能坍缩成一个闭式的一步跳转（以及这个闭式解用掉了哪三个条件），噪声调度如何决定训练样本落在信息量最大的那一段，以及为什么终点必然是高斯——那是抛物型方程的性质，与数据分布无关。
\item \hyperref[sec:14-Deep-Learning/08-Diffusion-Models/02]{``训练目标其实是一个去噪回归''}：模型学什么。变分下界如何拆成一串两个高斯之间的 KL、如何塌成平方距离、如何重参数化成``预测噪声''，以及最后丢掉时刻权重那一步为什么\emph{不是}推导而是一次经验选择。还有那个把整章串起来的等价：去噪与估计得分是同一件事。
\item \hyperref[sec:14-Deep-Learning/08-Diffusion-Models/03]{``反向采样是在拿步数换质量''}：怎么变回样本。确定性轨迹为什么允许跳步，随机轨迹为什么不允许，以及采样其实是在解一个概率流 ODE——步数就是积分步长，高阶求解器省步数与 Runge--Kutta 优于 Euler 是同一件事，低噪声端的刚性则给步长压缩画了一条线。
\item \hyperref[sec:14-Deep-Learning/08-Diffusion-Models/04]{``条件生成把控制信号写进得分''}：怎么让它听指令。条件得分等于无条件得分加一个分类器梯度，这一行 Bayes 推导是全篇的支点；classifier-free guidance 如何用同一个网络的两次预测替掉分类器，以及把强度调到 \(1\) 以上时，你采样的对象已经悄悄换成了一个被锐化的分布。
\end{itemize}

三条不要弄混的界线：``均匀加权的目标''不是变分下界，``跳步采样''不是对随机链的加速，``guidance 强度大于一''不是在更好地逼近条件分布。三处都是拿理论上的正确性换实践上的效果，换得值不值是经验问题，但换了什么必须清楚。

读完这一章，扩散模型的位置也就定下来了：它把生成的难度从训练搬到了采样。这个选择让它好训——损失一行、无对抗、收敛稳；也让它贵——每张图都要付几十次前向的成本。拆解不消灭困难，只改变困难出现的地方，而选在哪里付账，恰恰是这一类模型之间最实质的分歧。
